\subsection{Integral crystalline Lefschetz fails already on threefolds}

\begin{lemma}[A first-degree comparison of Hodge, de Rham, and crystalline dimensions]
\label{lem:OC-first-degree-torsion}
Let $V$ be a smooth projective geometrically connected scheme over a perfect field $k$ of characteristic $p>0$.  Put
\[
\begin{aligned}
 q(V)&=h^1(V,\OO_V),& h^{1,0}(V)&=h^0(V,\Omega_V^1),\\
 b(V)&=\rank_W H^1_{\cris}(V/W),&
 t(V)&=\dim_k H^2_{\cris}(V/W)[p].
\end{aligned}
\]
Then $H^1_{\cris}(V/W)$ is free and
\begin{equation}\label{eq:OC-first-degree-bounds}
 q(V)-b(V)\ \leq\ t(V)\ \leq\ q(V)+h^{1,0}(V)-b(V).
\end{equation}
\end{lemma}
\begin{proof}
Derived crystalline base change and the universal coefficient sequence give
\[
 0\to H^i_{\cris}(V/W)/p\to H^i_{\dR}(V/k)
       \to H^{i+1}_{\cris}(V/W)[p]\to0.
\]
In degree zero, the first arrow is the isomorphism $k\to k$, so $H^1_{\cris}[p]=0$.  Finiteness over the discrete valuation ring $W=W(k)$ proves freeness.  In degree one, consequently,
\[
                  \dim_k H^1_{\dR}(V/k)=b(V)+t(V).
\]
The Hodge spectral sequence bounds the left side above by $q(V)+h^{1,0}(V)$.  The conjugate spectral sequence bounds it below by $q(V)$: Cartier's isomorphism identifies its terms with
$H^a(V^{(1)},\Omega_{V^{(1)}}^b)$, and its term $(a,b)=(1,0)$ survives with no incoming or outgoing differential.  Equivalently, its low-degree edge sequence begins
\[
 0\longrightarrow H^1(V^{(1)},\OO_{V^{(1)}})
                 \longrightarrow H^1_{\dR}(V/k).
\]
For exact original locators, Cartier's isomorphism is
\cite[0, Theorem~2.1.9, p.~515]{Illusie1979DRW}; the displayed conjugate edge injection is explicitly written in
\cite[II, Section~6.14, p.~649]{Illusie1979DRW}.
The two bounds follow.  No degeneration of either spectral sequence is assumed.
\end{proof}

\begin{proposition}[A surface kernel produces integral threefold counterexamples]
\label{prop:OC-product-integral-counterexample}
\status{Proved; formal argument independently audited by the root and reduction-descent agents.}
Let $X$ be a smooth projective connected surface over an algebraically closed field of characteristic $p$, and let $A$ be very ample with
\[
                    h^1(X,A^{-1})=a>0.
\]
Choose $m\geq1$ satisfying
\[
                       a(m-1)>h^0(X,\Omega_X^1).
\]
Embed $Z=X\times\mathbf P^1$ by the very ample bundle
$B=A\boxtimes\OO_{\mathbf P^1}(m)$.
For every smooth hyperplane section $Y\in|B|$, both
\[
 T^2(Z)\longrightarrow T^2(Y),\qquad
 H^2_{\cris}(Z/W)\longrightarrow H^2_{\cris}(Y/W)
\]
fail, respectively, to be surjective on torsion and to have torsion-free cokernel.  More precisely,
\begin{equation}\label{eq:OC-product-torsion-growth}
 \dim_k T^2(Y)[p]-\dim_k T^2(Z)[p]
    \ \geq\ a(m-1)-h^0(X,\Omega_X^1)>0.
\end{equation}
Thus A.1(a) and A.1(b), as assertions for every smooth degree-one section, are false in dimension three.
\end{proposition}
\begin{proof}
Every smooth $Y$ is connected, since it is an ample effective divisor in a smooth connected projective threefold.  K\"unneth gives
\[
 H^1(Z,B^{-1})=0,
 \qquad
 H^2(Z,B^{-1})=
 H^1(X,A^{-1})\otimes H^1(\mathbf P^1,\OO(-m)).
\]
The second group has dimension $a(m-1)$.
The multiplication map
\[
                        H^2(Z,B^{-1})\longrightarrow H^2(Z,\OO_Z)
\]
is zero for every defining section of $Y$.  Indeed, write that section as a finite sum of tensors in
$H^0(X,A)\otimes H^0(\mathbf P^1,\OO(m))$.
Each tensor acts on the second K\"unneth factor through
$H^1(\mathbf P^1,\OO(-m))\to H^1(\mathbf P^1,\OO)=0$.
The divisor sequence therefore gives the exact sequence
\[
 0\longrightarrow H^1(X,\OO_X)\longrightarrow H^1(Y,\OO_Y)
   \longrightarrow H^1(X,A^{-1})\otimes H^1(\mathbf P^1,\OO(-m))
   \longrightarrow0.
\]
In particular,
\[
 q(Z)=q(X),\qquad q(Y)=q(X)+a(m-1),\qquad
 h^{1,0}(Z)=h^{1,0}(X).
\]
Rational crystalline weak Lefschetz makes the degree-one map an isomorphism and the degree-two map injective; the exact original input is \cite[Appendix, Section~9, Theorem~9.1 and Corollary~9.3, pp.~63--64]{NakkajimaYobuko2022} in version~3 dated 12~November~2022, whose proof is read in this run.  These are the current version numbers, not the Section~11 numbering of an earlier version.  Write
$b=b(Z)=b(Y)$; equivalently this is the common first Betti number, since integral crystalline $H^1$ is the covariant Dieudonn\'e realization of the reduced Picard variety
\cite[II, Proposition~3.11 and Remark~3.11.2, pp.~618--619]{Illusie1979DRW}.
Applying Lemma~\ref{lem:OC-first-degree-torsion} to $Y$ and $Z$ gives
\[
 t(Y)\geq q(X)+a(m-1)-b,
 \qquad t(Z)\leq q(X)+h^{1,0}(X)-b,
\]
which proves \eqref{eq:OC-product-torsion-growth}.

For a finite torsion module over $W$, the dimensions of its subgroup killed by $p$ and its quotient by $p$ agree; this number is its minimum number of generators.  A quotient cannot have more generators than its source.  Hence $T^2(Z)\to T^2(Y)$ is not surjective.
Let $u:H^2_{\cris}(Z/W)\to H^2_{\cris}(Y/W)$.
Its rational injectivity implies
\[
            \im(u)\cap T^2(Y)=u\bigl(T^2(Z)\bigr).
\]
Indeed, a source element whose image is torsion becomes zero after tensoring with $W[1/p]$, so is itself torsion.
Thus the nonzero finite torsion module
$T^2(Y)/u(T^2(Z))$ injects into $\coker(u)$.
The latter is not torsion-free.  This argument does not claim that $u$ has a nonzero kernel; A.1(c) is a separate issue.
\end{proof}

\begin{lemma}[A numerical differential bound for the explicit surface]
\label{lem:OC-explicit-differential-bound}
\status{Proved; independently audited by the root and reduction-descent agents.}
For the characteristic-two surface $X$ and hyperplane bundle $A$ of
Corollary~\ref{cor:OC-surface-exact-degrees},
\[
                          h^0(X,\Omega_X^1)\leq591\cdot9^{556}.
\]
\end{lemma}
\begin{proof}
Retain the maps $X\xrightarrow{\rho}S\xrightarrow{\psi}P\xrightarrow{\pi}C$ from that construction.  Here
$\psi$ is the tame cubic cover of $P=\mathbf P_C(E)$, branched at the disjoint smooth curves $E_0$ and $G$; $g(C)=7$, $\deg N=2$, $\det E=N^3$, and
$M=\OO_P(-1)\otimes\pi^*N^2$.
A tame root cover with simple normal-crossing branch divisor $D$ has
\[
 \Omega^1_{\rm cover}\ \lhook\joinrel\longrightarrow
                \text{pullback of }\Omega^1_{\rm base}(\log D).
\]
This can be checked at every point directly: if a branch parameter $s$ is replaced by $u$ with $u^q=s$, then $ds/s=q\,du/u$ and $q$ is invertible.  At a double intersection use the two branch parameters; away from the branch the cover is \'etale.  Thus the asserted inclusion is an inclusion of locally free sheaves and needs no statement about wild logarithmic covers.

Apply it first to $\psi$.  Projection formula gives an upper bound for $h^0(S,\Omega_S^1)$ by the sum of the three dimensions
\[
 H^0\bigl(P,\Omega_P^1(\log(E_0+G))\otimes M^i\bigr),
                              \qquad i=0,1,2.
\]
The relative cotangent sequence on the projective line bundle shows
\[
 h^0(P,\Omega_P^1)=7,
 \qquad H^0(P,\Omega_P^1\otimes M^i)=0\quad(i=1,2).
\]
For the latter assertion, its two end terms restrict to
$\OO_{\mathbf P^1}(-i)$ and $\OO_{\mathbf P^1}(-2-i)$ on every fibre; both have no sections.  For $i=0$ the first end term is $\pi^*\Omega_C^1$ and the second still has negative fibre degree.
The residue exact sequence, tensored with $M^i$, now bounds these logarithmic dimensions.
For $i=0$ the two residues contribute at most two, so the bound is $7+2=9$.
For $i=1,2$ one has
\[
 M^i|_{E_0}=N^{-i},\qquad
 M^i|_G=(\pi|_G)^*N^{2i}.
\]
The first has no section.  The second has degree $8i$, since $G\to C$ has degree two; here $\OO_P(1)|_G$ is trivial because $G$ is disjoint from the divisor $E_0$.  On any smooth connected projective curve a line bundle of nonnegative degree $e$ has at most $e+1$ sections, by successive subtraction of points.  The remaining two bounds are therefore $9$ and $17$.  Altogether,
\[
                         h^0(S,\Omega_S^1)\leq9+9+17=35.
\]

The cover $\rho$ has the $556$ smooth irreducible branch curves $D_i$, with pairwise transverse intersections and no triple intersection.  Its degree is $9^{556}$, and
\[
 \rho_*\OO_X
    =\bigoplus_{0\leq a_i<9}L^{-\sum a_i}.
\]
The same local logarithmic inclusion and projection formula bound $h^0(X,\Omega_X^1)$ by the sum of
\[
 h^0\!\left(S,\Omega_S^1\left(\log\sum D_i\right)
                              \otimes L^{-\sum a_i}\right).
\]
Each summand is at most $h^0(S,\Omega_S^1(\log\sum D_i))$: multiplication by a fixed nonzero global section of $L^{\sum a_i}$ gives an injection into that sheaf, since $S$ is integral and the logarithmic sheaf is locally free.  Finally, the residue sequence bounds this last dimension by
$h^0(S,\Omega_S^1)+556\leq591$.
There are exactly $9^{556}$ summands, proving the stated numerical bound.
\end{proof}

\begin{theorem}[An explicit degree-one integral counterexample]
\label{thm:OC-explicit-integral-counterexample}
\status{Proved; independently audited by the root and reduction-descent agents.}
In characteristic two, let $(X,A)$ be the surface in
Corollary~\ref{cor:OC-surface-exact-degrees} and set
\[
             m=591\cdot9^{556}+2,\qquad
             Z=X\times\mathbf P^1,\qquad B=A\boxtimes\OO(m).
\]
The displayed surface embedding and the degree-$m$ Veronese embedding of $\mathbf P^1$, followed by the Segre embedding, realize $B$ as a hyperplane bundle on a smooth connected projective threefold.  For every smooth hyperplane section $Y$ of this threefold,
\[
 \dim_k H^2_{\cris}(Y/W)[2]
       \geq \dim_k H^2_{\cris}(Z/W)[2]+1.
\]
In particular its crystalline vanishing quotient has nonzero $2$-power torsion, and restriction does not identify the degree-two torsion subgroups.
\end{theorem}
\begin{proof}
Here $a=h^1(X,A^{-1})=1$.  Lemma~\ref{lem:OC-explicit-differential-bound} gives
$m-1>h^0(X,\Omega_X^1)$, so all assertions follow from
Proposition~\ref{prop:OC-product-integral-counterexample}.
Smooth sections exist by Bertini; an independent generic linear combination of the displayed Segre sections gives a completely specified example over the algebraic closure of the resulting rational function field.  All exponents and dimensions entering the construction are explicit integers.
\end{proof}


\begin{corollary}[Explicit integral counterexamples in every prime]
\label{cor:OC-all-prime-integral-counterexample}
\status{Proved; independently audited by the root and reduction-descent agents.}
For every prime $p$, take the surface $(S,L)$ of
Proposition~\ref{prop:OC-generated-Raynaud}.  Set
\[
 \ell=p+1,\quad q=\ell^2,\quad
 n=\deg N=\begin{cases}2,&p=2,\\p(p-2),&p\geq3,\end{cases}
 \qquad g=1+\frac{p\ell n}{2},
\]
\[
 B_p=g+p+2+\frac{p^3(p+1)n}{2},
 \qquad a_i=p\ell-i\quad(0\leq i\leq p),
\]
and define the positive integer
\[
 r=\sum_{i=0}^p\left[
 (a_i+1)\bigl((q+pi)n+1-g\bigr)
     +\frac{\ell n\,a_i(a_i+1)}2\right].
\]
The $q$th-root cover with one root for every member of a generic basis of
$H^0(S,L^q)$ has degree $q^r$ and produces a very ample pair $(X,A)$.
Put
\[
           m=q^r(B_p+r)+2,\qquad Z=X\times\mathbf P^1,
           \qquad B=A\boxtimes\OO(m).
\]
For every smooth degree-one hypersurface section $Y$ in this embedding,
$H^2_{\cris}(Y/W)[p]$ has strictly larger dimension than
$H^2_{\cris}(Z/W)[p]$.
Thus A.1(a) and A.1(b) fail in dimension three for every prime.
\end{corollary}
\begin{proof}
The proof of Lemma~\ref{lem:OC-explicit-differential-bound} works for the tame degree-$p+1$ surface cover.  The untwisted logarithmic differential term contributes at most $g+2$.  For every $i=1,\ldots,p$, the plain differential twist has no sections by its negative relative fibre degrees, the residue at the ramification section has negative degree, and the residue on $G$ has degree $p^2in$.  This last identity uses $\deg(G/C)=p$ and $M^i|_G=(\pi|_G)^*N^{pi}$.  Hence
\[
 h^0(S,\Omega_S^1)
 \leq g+2+\sum_{i=1}^p(p^2in+1)=B_p.
\]
For the dimension $r$, use $L^q=\psi^*F$ with
$F=\OO_P(p\ell)\otimes\pi^*N^q$.
Its $i$th eigensheaf pushes to
$\operatorname{Sym}^{a_i}E\otimes N^{q+pi}$ on $C$, whose filtered line quotients are $N^{q+pi+\ell j}$ for $0\leq j\leq a_i$.
Their degrees exceed $2g-2=p\ell n$, so their $H^1$ groups vanish.  Riemann--Roch and addition along the filtration give exactly the displayed formula for $r=h^0(S,L^q)$.
The multiroot logarithmic bound is consequently
\[
                         h^0(X,\Omega_X^1)\leq q^r(B_p+r).
\]
Meanwhile $h^1(X,A^{-1})\geq h^1(S,L^{-1})\geq1$ by the direct summand in Lemma~\ref{lem:OC-Kummer-very-ample}.
Our choice of $m$ satisfies the strict inequality in
Proposition~\ref{prop:OC-product-integral-counterexample}, which proves the assertion.
At $p=2$ the formula gives $B_2=35$, $r=556$, and recovers exactly the preceding numerical example.
\end{proof}

\begin{remark}[Exact degree refinement and the odd-prime additional examples]
\label{rem:OC-integral-degree-gate}
The exact degree range for the characteristic-two counterexample is established in Theorem~\ref{thm:OC-product-exact-degrees}: A.1(a) and A.1(b) fail in degree one and hold in every degree at least two, while integral restriction is injective in every degree.  Corollary~\ref{cor:OC-all-prime-product-degrees} gives the explicit bound $d\geq p^2-p$ for the additional all-prime examples; their intervening odd-prime degree refinements remain separate.  The proof uses the relative de Rham differential and both universal-coefficient diagrams.  Coherent equality alone, or an isomorphism only on subgroups killed by $p$, would not imply the full torsion assertion: multiplication by $p$ gives $W/p\hookrightarrow W/p^2$ and illustrates the latter insufficiency.
\end{remark}
